
\documentclass[10pt,aspectratio=169]{beamer}

% Sử dụng theme màu xanh nhẹ nhàng
\usetheme{Madrid}
\definecolor{ptitblue}{RGB}{40,90,170}
\definecolor{ptitblue2}{RGB}{100,160,220}
\definecolor{lightblue}{RGB}{225,236,254}
\definecolor{tableblue}{RGB}{230,242,255}

% Các tông xanh chủ đạo
\setbeamercolor{structure}{fg=ptitblue}
\setbeamercolor{frametitle}{fg=white,bg=ptitblue}
\setbeamercolor{title}{fg=white,bg=ptitblue}
\setbeamercolor{block title}{fg=white,bg=ptitblue}
\setbeamercolor{block body}{fg=black,bg=lightblue}
\setbeamercolor{footer}{fg=white,bg=ptitblue}

\usepackage[utf8]{inputenc}
\usepackage[T5]{fontenc}
\usepackage[vietnamese]{babel}
\usepackage{amsmath}
\usepackage{booktabs}
\usepackage{algorithm}
\usepackage{algpseudocode}
\usepackage{tcolorbox}
\usepackage{array}
\usepackage{tikz}
\usepackage{colortbl}
\usepackage{xcolor}

% --- Patch for Đ/đ in T5 Vietnamese font ---
\usepackage{newunicodechar}
\newunicodechar{Đ}{\DJ}
\newunicodechar{đ}{\dj}
\DeclareUnicodeCharacter{0110}{\DJ}
\DeclareUnicodeCharacter{0111}{\dj}

% --- Patch cho table alternate row colors với xanh nhạt ---
\definecolor{rowblue1}{RGB}{230,242,255}
\definecolor{rowblue2}{RGB}{255,255,255}
\newcommand{\rowcolorsalt}{\rowcolors{2}{rowblue1}{rowblue2}}

% --- Custom footer: chỉ hiện 'Nhóm 1' bên trái và đếm số trang (slide) bên phải ---
\setbeamertemplate{footline}{
  \leavevmode%
  \hbox{%
    \begin{beamercolorbox}[wd=.5\paperwidth,ht=2.5ex,dp=1.0ex,left]{footer}
      \hspace{1.5em}Nhóm 1
    \end{beamercolorbox}
    \begin{beamercolorbox}[wd=.5\paperwidth,ht=2.5ex,dp=1.0ex,right]{footer}
      \insertframenumber{} / \inserttotalframenumber\hspace{1.5em}
    \end{beamercolorbox}
  }%
}

\setbeamertemplate{navigation symbols}{}

\title[Đồ án IR]{Xây dựng hệ thống truy xuất thông tin trên tập dữ liệu SciFact}
\author{Trần Đức Trình -- B22DCKH122 \and Nguyễn Ngọc Quang -- B22DCKH092}
\institute{Học viện Công nghệ Bưu chính Viễn thông}
\date{\the\year}

\begin{document}

% --- Trang bìa đơn giản, chỉ tên đề tài (nhỏ) + thành viên + giảng viên ---
\begin{frame}[plain]
    \centering

    % Tên đề tài (nhỏ hơn, đơn giản)
    {\bfseries\color{ptitblue}\Large Xây dựng hệ thống truy xuất thông tin trên tập dữ liệu SciFact}

    \vspace{1.5cm}

    % Giảng viên hướng dẫn
    {\normalsize \textbf{GVHD:} ThS. Phạm Đức Cường}

    \vspace{0.7cm}

    % Thành viên nhóm
    {\normalsize 
        \textbf{Nhóm 1:}\\[2pt]
        Trần Đức Trình -- B22DCKH122\\
        Nguyễn Ngọc Quang -- B22DCKH092
    }
\end{frame}

\begin{frame}{Tổng quan}
\begin{tcolorbox}[colback=lightblue,colframe=ptitblue,title=Vấn đề cần giải quyết]
\smallskip
\textbf{Bài toán}: Tìm kiếm và xếp hạng tập tài liệu $D$ theo mức độ liên quan với truy vấn $Q$.
\end{tcolorbox}

\vspace{.7em}
\begin{columns}[T,onlytextwidth]
\column{0.5\textwidth}
\begin{itemize}
\item 5.183 tài liệu khoa học y sinh (SciFact)
\item 300 truy vấn dạng phát biểu khoa học
\end{itemize}

\column{0.5\textwidth}
\begin{itemize}
\item So sánh: \textbf{BM25}, \textbf{BM25F}, \textbf{PRF-Rocchio}
\item Các chỉ số: Precision@K, Recall@K, MAP
\end{itemize}
\end{columns}
\end{frame}

\begin{frame}{Pipeline}
\vspace{-0.5em}
\centering
\begin{tikzpicture}[node distance=2.2cm,>=latex]
\tikzstyle{stage}=[rectangle,draw=ptitblue,fill=lightblue,rounded corners=3pt,minimum height=1.03cm,minimum width=3.3cm,align=center,font=\normalsize]
\node[stage] (raw) {Raw Data};
\node[stage,right of=raw,xshift=1.5cm] (pre) {Preprocessing};
\node[stage,right of=pre,xshift=1.5cm] (idx) {Indexing \& Ranking};
\node[stage,right of=idx,xshift=1.5cm] (eva) {Evaluation};

\draw[->,thick, ptitblue] (raw) -- (pre);
\draw[->,thick, ptitblue] (pre) -- (idx);
\draw[->,thick, ptitblue] (idx) -- (eva);
\end{tikzpicture}

\vspace{1em}
\begin{columns}[T,onlytextwidth]
\column{0.5\textwidth}
\begin{itemize}
    \item Tokenize, stopword removal, stemming
    \item Inverted Index, Postings list
\end{itemize}
\column{0.5\textwidth}
\begin{itemize}
    \item BM25, BM25F, PRF-Rocchio
    \item Evaluation: Precision@K, Recall@K, MAP
\end{itemize}
\end{columns}
\end{frame}

% ---------- BM25 Theory Slide (Slide 4) ----------
\begin{frame}{BM25}
\begin{tcolorbox}[colback=lightblue,boxrule=0pt]
\textbf{Giới thiệu:}
BM25 là phương pháp xếp hạng dựa trên mô hình xác suất cho truy xuất văn bản. Ý tưởng chính là đánh giá mức độ liên quan giữa tài liệu $D$ và truy vấn $Q$ thông qua tần suất xuất hiện từ truy vấn trong tài liệu, đồng thời chuẩn hóa độ dài tài liệu.
\end{tcolorbox}

\vspace{0.7em}

\textbf{Công thức tính điểm:}
\[
\text{score}(D, Q) = \sum_{t \in Q}
\log\left(\frac{N-df_t+0.5}{df_t+0.5}\right) \cdot
\frac{tf_{t,D} \cdot (k_1 + 1)}{tf_{t,D} + k_1 \cdot \left[1-b+b\frac{L_D}{L_{ave}}\right]}
\]
\begin{itemize}
    \item $tf_{t,D}$: Tần suất từ $t$ trong tài liệu $D$.
    \item $N$: Tổng số tài liệu, $df_t$: số tài liệu chứa $t$.
    \item $k_1$: Tham số bão hòa TF (thường 1.2--2.0), $b$: Tham số chuẩn hóa độ dài (thường 0.75).
    \item $L_D$, $L_{ave}$: Độ dài tài liệu $D$ và độ dài trung bình.
\end{itemize}
\end{frame}

\begin{frame}{BM25 Pseudocode}
\vspace{-0.7em}
\scriptsize
\begin{tcolorbox}[colback=lightblue,colframe=ptitblue,boxrule=0.4pt,title=BM25 Algorithm (Pseudocode)]
\vspace{-1em}
\begin{algorithm}[H]
\begin{algorithmic}[1]
\State Initialize \texttt{Scores} as an empty dictionary
\For{each term $t$ in query $Q$}
    \If{$t$ exists in index}
        \State Retrieve postings and $df_t$
        \State Compute $IDF_t = \log \left(\frac{N-df_t+0.5}{df_t+0.5}\right)$
        \For{each document $d$ in postings}
            \State $tf =$ term frequency of $t$ in $d$
            \State $norm = 1 - b + b \frac{L_d}{L_{ave}}$
            \State $score = IDF_t \times \frac{tf \times (k_1 + 1)}{tf + k_1 \cdot norm}$
            \State Add $score$ to \texttt{Scores[$d$]}
        \EndFor
    \EndIf
\EndFor
\State \Return \texttt{Scores}
\end{algorithmic}
\end{algorithm}
\end{tcolorbox}
\normalsize
\end{frame}

% ---------- BM25F Theory Slide (Slide 6) ----------
\begin{frame}{BM25F}
\begin{tcolorbox}[colback=lightblue,boxrule=0pt]
\textbf{Giới thiệu:}
BM25F là phiên bản mở rộng của BM25, áp dụng cho tài liệu có nhiều trường, ví dụ như ``title'', ``abstract''. Thuật toán gán trọng số cho từng trường để tận dụng thông tin cấu trúc của tài liệu.
\end{tcolorbox}

\vspace{0.6em}
\textbf{Công thức:}
\[
\tilde{tf}_{t,D} = \sum_{f} w_f \cdot \frac{tf_{t,D,f}}{1-b_f+b_f\frac{L_{D,f}}{L_{ave,f}}}
\]
\[
\text{score}(D, Q) = \sum_{t\in Q}
\log\left(\frac{N-df_t+0.5}{df_t+0.5}\right) \cdot
\frac{\tilde{tf}_{t,D} \cdot (k_1+1)}{\tilde{tf}_{t,D}+k_1}
\]
\vspace{-0.5em}
\begin{itemize}
    \item $w_f$: Trọng số của field $f$ (ví dụ $w_{title}$, $w_{abstract}$).
    \item $tf_{t,D,f}$: Tần suất $t$ ở field $f$ của tài liệu $D$.
    \item $b_f$: Tham số chuẩn hóa độ dài cho từng field.
\end{itemize}
\end{frame}

\begin{frame}{BM25F Pseudocode}
\scriptsize
\begin{tcolorbox}[colback=lightblue,colframe=ptitblue,boxrule=0.4pt,title=BM25F Algorithm (Pseudocode)]
\vspace{-1em}
\begin{algorithm}[H]
\begin{algorithmic}[1]
\State Initialize \texttt{Scores} as empty map
\For{each term $t$ in $Q$}
    \If{$t$ in index}
        \State Compute $IDF_t$
        \For{each document $d$ containing $t$}
            \State $\tilde{tf} \leftarrow$ 0
            \For{each field $f$}
                \State $\tilde{tf} \mathrel{+}= w_f \cdot$ normalized tf of $t$ in $f$ of $d$
            \EndFor
            \State $score \leftarrow IDF_t \times \frac{\tilde{tf} (k_1+1)}{\tilde{tf} + k_1}$
            \State Add $score$ to \texttt{Scores[$d$]}
        \EndFor
    \EndIf
\EndFor
\State \Return \texttt{Scores}
\end{algorithmic}
\end{algorithm}
\end{tcolorbox}
\normalsize
\end{frame}

% ---------- PRF-Rocchio Theory Slide (Slide 8) ----------
\begin{frame}{PRF-Rocchio}
\begin{tcolorbox}[colback=lightblue,boxrule=0pt]
\textbf{Giới thiệu:}
PRF (Pseudo Relevance Feedback) với Rocchio là phương pháp mở rộng truy vấn bằng cách tận dụng các tài liệu top-$K$ được trả về từ kết quả ban đầu, giả sử chúng là tài liệu liên quan. Ý tưởng: truy vấn mới = truy vấn gốc + thông tin lấy từ tài liệu liên quan.
\end{tcolorbox}

\vspace{0.7em}
\textbf{Công thức Rocchio:}
\[
\vec{q}_m = \alpha\vec{q}_0 + \frac{\beta}{|D_r|}\sum_{\vec{d}\in D_r}\vec{d}
\]
\begin{itemize}
    \item $\vec{q}_0$: Truy vấn gốc dạng vector.
    \item $D_r$: Tập các tài liệu top-$K$.
    \item $\alpha$: Trọng số giữ độ trung thành truy vấn gốc.
    \item $\beta$: Trọng số mở rộng.
\end{itemize}
\vspace{0.3em}
\textbf{Lưu ý:} Nếu top-$K$ chứa nhiễu, có thể gây \textit{Query Drift.}
\end{frame}

\begin{frame}{PRF-Rocchio Pseudocode}
\scriptsize
\begin{tcolorbox}[colback=lightblue,colframe=ptitblue,boxrule=0.4pt,title=PRF-Rocchio Algorithm (Pseudocode)]
\vspace{-1em}
\begin{algorithm}[H]
\begin{algorithmic}[1]
\State $R \leftarrow$ top-$K$ docs from $D$ using $q_0$
\State $centroid \leftarrow$ mean of $R$
\State $q_m \leftarrow \alpha q_0 + \beta\, centroid$
\State \Return Retrieve results from $D$ using $q_m$
\end{algorithmic}
\end{algorithm}
\end{tcolorbox}
\normalsize
\end{frame}

\begin{frame}{Xử lý và Indexing}
\vspace{-1em}
\begin{columns}[T,onlytextwidth]
\column{0.54\textwidth}
\begin{block}{Tiền xử lý}
\begin{itemize}
    \item Regex Tokenize $\rightarrow$ lowercase $\rightarrow$ stopword removal $\rightarrow$ Porter Stemming
\end{itemize}
\end{block}

\column{0.44\textwidth}
\begin{block}{Indexing}
\begin{itemize}
    \item \textbf{Inverted Index}: $Index[\text{term}] \rightarrow$ Postings list $(\text{doc\_id}, tf)$
    \item \textbf{BM25F}: $Index\_F[\text{term}] \rightarrow$ (doc\_id, tf dictionary theo field)
\end{itemize}
\end{block}
\end{columns}

\vspace{0.7em}
\vspace{-0.2em}
\centering
\footnotesize
\colorbox{lightblue}{Tập dữ liệu: 5.183 documents, 300 queries (SciFact)}
\end{frame}

\begin{frame}{Tham số và chỉ số đánh giá}
\begin{columns}[T,onlytextwidth]
\column{0.49\textwidth}
\begin{block}{Thông số}
\begin{itemize}
    \item $k_1=1.5$, $b=0.75$ (BM25)
    \item $w_{title}=2.0,\  w_{text}=1.0$, $b_{title}=0.8$, $b_{text}=0.75$ (BM25F)
    \item $\alpha=1.0,\, \beta=0.5,\, K=1, \, T=3$ (PRF)
\end{itemize}
\end{block}

\column{0.49\textwidth}
\begin{block}{Chỉ số}
\begin{itemize}
    \item Precision@K: Độ chính xác top-$K$
    \item Recall@K: Tỷ lệ hồi tưởng top-$K$
    \item MAP: Trung bình AP trên toàn bộ truy vấn
\end{itemize}
\end{block}
\end{columns}
\end{frame}

\begin{frame}{Kết quả thực nghiệm}
\begin{table}
\centering
\renewcommand{\arraystretch}{1.5}
\setlength{\tabcolsep}{15pt}
\large
\rowcolors{2}{rowblue1}{rowblue2}
\begin{tabular}{>{\bfseries}l c c c c}
\toprule
Phương pháp & P@1 & P@10 & Recall@10 & MAP \\
\midrule
BM25 Baseline & 0.5600 & 0.0907 & 0.8202 & 0.6480 \\
BM25F & 0.5633 & 0.0907 & 0.8202 & 0.6524 \\
BM25 + PRF & 0.5600 & 0.0920 & 0.8277 & 0.6512 \\
\bottomrule
\end{tabular}
\end{table}

\vspace{.6em}
\begin{itemize}
    \item \textbf{BM25F} cải thiện toàn diện và đạt kết quả tốt nhất (MAP 0.6524), chứng minh giá trị của việc tận dụng cấu trúc title-abstract để xếp hạng.
    \item \textbf{PRF} tối ưu hóa bằng TF-IDF và bộ lọc stopword giúp tránh hiện tượng Topic Drift, qua đó cải thiện rõ rệt Recall và MAP so với Baseline một cách an toàn.
\end{itemize}
\end{frame}

\begin{frame}{Kết luận}
\begin{block}{Tóm tắt}
\begin{itemize}
    \item \textbf{BM25F} lập mốc độ chính xác tổng thể (MAP, P@1) cao nhất, khẳng định vai trò qua trọng của việc phân rã trường dữ liệu.
    \item \textbf{PRF} nâng cao khả năng thu hồi (Recall) thông qua thuật toán mở rộng truy vấn có chọn lọc rủi ro, hoàn thiện hệ thống từ hướng xử lý câu hỏi.
\end{itemize}
\end{block}
\vspace{1em}
\begin{block}{Hướng cải thiện}
\begin{itemize}
    \item Kết hợp các mô hình biểu diễn ngữ nghĩa hiện đại như BERT để cải thiện chất lượng truy vấn và kết quả truy xuất.
    \item Ứng dụng học sâu hoặc mô hình rerank để tăng khả năng nhận diện văn bản liên quan.
    \item Triển khai hệ thống trên nhiều tập dữ liệu khác nhau.
\end{itemize}
\end{block}
\end{frame}

\end{document}